\begin{textsample}{POS Dim 2 – ms – Score 22.00 – ms/ms\_em\_39.txt}  \label{ex:f2_pos_045}
4.4.1.
Distribuição da \textbf{Carga} \textbf{Horária} e Organização da \textbf{Oferta} Segundo a LDB, a \textbf{carga} \textbf{horária} mínima anual para o Novo Ensino Médio será de mil horas, devendo ser ampliada, de forma progressiva, para mil e quatrocentas horas.
Este \textbf{Currículo} de Referência sugere o desenvolvimento da Formação Geral Básica ( FGB ) em 600 ( seiscentas ) horas anuais e dos \textbf{Itinerários} \textbf{Formativos} ( IF ) em 400 ( quatrocentas ) horas anuais.
Assim, semanalmente, serão \textbf{ofertadas} 30 ( trinta ) horas-aulas, divididas em 18 ( dezoito ) para a FGB e 12 ( doze ) para os IF.
Diante desse arranjo, é possível identificar algumas possibilidades de organização da \textbf{carga} \textbf{horária} \textbf{semanal} : As instituições de ensino poderão propor formas de distribuição \textbf{semanal} da \textbf{carga} \textbf{horária}, de modo a possibilitar o desenvolvimento da Formação Geral Básica e do Itinerário \textbf{Formativo} em todos os dias da semana ( Matriz 1 ) ou em dias da semana específicos ( Matriz 2 ). 4.5. \textbf{Itinerário} \textbf{Formativo} Propedêutico Os \textbf{Itinerários} \textbf{Formativos} são definidos pelas \textbf{Diretrizes} Curriculares Nacionais para o Ensino Médio, conforme se transcreve : > Art. 6º, III - \textbf{itinerários} \textbf{formativos} : cada conjunto de unidades curriculares \textbf{ofertadas} pelas escolas e \textbf{redes} de ensino que possibilitam ao estudante aprofundar seus conhecimentos e se preparar para o \textbf{prosseguimento} de estudos ou para o mundo do trabalho de forma a contribuir para a construção de soluções de problemas específicos da sociedade. ( RESOLUÇÃO CNE/CEB n. 3/2018, Art. 6º, Inciso III ).
Em sua composição, os IF compreendem um conjunto de situações e atividades educativas que os estudantes podem escolher, conforme seus interesses, para aprofundar e ampliar aprendizagens em uma ou mais Áreas de Conhecimento e/ou na FTP, segundo a Resolução CNE/CEB n. 3/2018 : I - Linguagens e suas Tecnologias : aprofundamento de conhecimentos \textbf{estruturantes} para aplicação de diferentes linguagens em contextos sociais e de trabalho, construindo arranjos curriculares que permitam estudos em línguas vernáculas, estrangeiras, clássicas e indígenas, Língua Brasileira de Sinais ( LIBRAS ), das artes, design, linguagens digitais, corporeidade, artes cênicas, roteiros, produções literárias, dentre outros, considerando o contexto local e as possibilidades de \textbf{oferta} pelos sistemas de ensino ; II - Matemática e suas Tecnologias : aprofundamento de conhecimentos \textbf{estruturantes} para aplicação de diferentes conceitos matemáticos em contextos sociais e de trabalho, construindo arranjos curriculares que permitam estudos em resolução de problemas e análises complexas, funcionais e não lineares, análise de dados estatísticos e probabilidade, geometria e topologia, robótica, automação, inteligência artificial, programação, jogos digitais, sistemas dinâmicos, dentre outros, considerando o contexto local e as possibilidades de \textbf{oferta} pelos sistemas de ensino ; III - Ciências da Natureza e suas Tecnologias : aprofundamento de conhecimentos \textbf{estruturantes} para aplicação de diferentes conceitos em contextos sociais e de trabalho, organizando arranjos curriculares que permitam estudos em astronomia, metrologia, física geral, clássica, molecular, quântica e mecânica, instrumentação, ótica, acústica, química dos produtos naturais, análise de fenômenos físicos e químicos, meteorologia e climatologia, microbiologia, imunologia e parasitologia, ecologia, nutrição, zoologia, dentre outros, considerando o contexto local e as possibilidades de \textbf{oferta} pelos sistemas de ensino ; IV - Ciências Humanas e Sociais Aplicadas : aprofundamento de conhecimentos \textbf{estruturantes} para aplicação de diferentes conceitos em contextos sociais e de trabalho, construindo arranjos curriculares que permitam estudos em relações sociais, modelos econômicos, processos políticos, pluralidade cultural, historicidade do universo, do homem e natureza, dentre outros, considerando o contexto local e as possibilidades de \textbf{oferta} pelos sistemas de ensino ; V - Formação \textbf{Técnica} e Profissional : desenvolvimento de programas educacionais inovadores e \textbf{atualizados} que promovam efetivamente a \textbf{qualificação} \textbf{profissional} dos estudantes para o mundo do trabalho, com vistas a sua \textbf{habilitação} \textbf{profissional} para o desenvolvimento de vida e \textbf{carreira} e adaptar -se às novas condições ocupacionais e às exigências do mundo do trabalho contemporâneo e suas contínuas transformações, em condições de competitividade, produtividade e inovação, considerando o contexto local e as possibilidades de \textbf{oferta} pelos sistemas de ensino ( BRASIL, 2018a ).
Contemplados na parte \textbf{flexível} do \textbf{Currículo}, os \textbf{Itinerários} \textbf{Formativos} fortalecem, ampliam e aprofundam o interesse do estudante pela área escolhida, incentivam a construção de um projeto de vida, o desenvolvimento do protagonismo juvenil, de habilidades e de valores, além do desenvolvimento de competências socioemocionais, como descritas na \textbf{Portaria} MEC n. 1.432/2018 ( BRASIL. 2018d ).
Os \textbf{Itinerários} \textbf{Formativos} Propedêuticos estão correlacionados às Áreas de Conhecimento, devem considerar as demandas e necessidades do mundo contemporâneo, estar \textbf{sintonizados} com os diferentes interesses dos estudantes e sua inserção na sociedade, o contexto local e as possibilidades de \textbf{oferta} dos sistemas e instituições de ensino.
Orientados para o aprofundamento e ampliação das aprendizagens em áreas do conhecimento, os \textbf{Itinerários} devem garantir a apropriação de procedimentos cognitivos e o uso de metodologias que favoreçam o protagonismo juvenil e organizarem -se em torno de um ou mais Eixos \textbf{Estruturantes}.
De forma a permitir a efetiva possibilidade de escolha e o protagonismo dos estudantes, as UCs devem promover o aprofundamento e a ampliação das aprendizagens essenciais de uma ou mais Áreas de Conhecimento, que consolidem e aprofundem a formação integral, a incorporação de valores universais, ampliem a visão de mundo dos estudantes e auxiliem a tomar decisões e agir com autonomia e responsabilidade.
As Áreas de Conhecimento devem proporcionar a apropriação de conceitos e categorias básicas que estabelecem um conjunto necessário de saberes integrados e significativos.
Assim, o Aprofundamento deve buscar expandir os aprendizados promovidos pela Formação Geral Básica, \textbf{viabilizar} a ampliação dos aprendizados em articulação com temáticas contemporâneas, sintonia com o contexto e os interesses dos estudantes, explorar potenciais e \textbf{vocações}, possibilitar um histórico escolar personalizado com maior tempo de dedicação a um Itinerário \textbf{Formativo} escolhido em consonância com o seu Projeto de Vida.

% matched lemmas: atualizar, carga, carreira, currículo, diretriz, estruturante, flexível, formativo, habilitação, horário, itinerário, oferta, ofertar, portaria, profissional, prosseguimento, qualificação, rede, semanal, sintonizar, técnico, viabilizar, vocação
\end{textsample}
